%!TEX root = ../template.tex
%% ====================================================================
%% Anexo G -- Definicoes matematicas dos instrumentos de avaliacao
%% Esqueleto criado 2026-09-08. Recebe a matematica que os Caps. 6 e 7
%% referem (Anexo~\ref{app:matematica}): estimadores, estatisticas de teste,
%% construcao das bandas, bootstrap de blocos, criterio de aceitacao,
%% referencia congelada, monitores e gerador da simulacao.
%% Por escrever.
%% ====================================================================

\chapter{Definições matemáticas dos instrumentos de avaliação}
% [REV-2026-09-13][AP-G][P2][APENDICE] A matemática essencial não pode permanecer num apêndice vazio
% Priorizar: modelo anual e sigma com graus de liberdade; seleção e
% desfasamentos civis; estimandos de perdas/cobertura; algoritmo MBB e
% condicionamento/reestimação; Wald HAC e famílias; Huber/Theil-Sen;
% centro e largura de cada banda; regra 4/5; referência/estado/inovação;
% recorrência CUSUM com k,h/reset; gerador, calibração por horizonte e
% censura de deteção. Manter no corpo as equações necessárias à leitura.
% Um pseudocódigo de um ano de treino/teste e outro do simulador ajudam
% mais que reproduzir funções completas do código.
% [/REV-2026-09-13]
\label{app:matematica}

%% >> Por secao do Cap. 6 e 7: notacao; regressao anual e replicacao;
%% >> skill e contraste emparelhado de perdas com IC MBB; Wald HAC empilhado
%% >> e familias BH/BY; cobertura e IC MBB; modelos encaixados e valor
%% >> incremental; estimadores de Huber e Theil--Sen; bandas AR e quantilica;
%% >> criterio de aceitacao; referencia congelada, escala e populacao
%% >> primaria; monitores (residuo bruto, filtro de nivel, soma acumulada,
%% >> canais de carga e escala); gerador Monte Carlo, calibracao de limiares
%% >> e desvio minimo detetavel.
